\documentclass[a4paper,11pt,oneside]{memoir}

% Castellano
\usepackage[spanish,es-tabla]{babel}
\selectlanguage{spanish}
\usepackage[utf8]{inputenc}
\usepackage{placeins}

\RequirePackage{booktabs}
\RequirePackage[table]{xcolor}
\RequirePackage{xtab}
\RequirePackage{multirow}

%Float
\usepackage{float}

%Simbolos
\usepackage{textcomp}


%Citas bibliografía
\usepackage{cite}
\usepackage[numbers,sort]{natbib}

% Links
\usepackage[colorlinks]{hyperref}
\hypersetup{
	allcolors = {blue}
}

% Ecuaciones
\usepackage{amsmath}

% Rutas de fichero / paquete
\newcommand{\ruta}[1]{{\sffamily #1}}

% Párrafos
\nonzeroparskip


% Imagenes
\usepackage{graphicx}
\newcommand{\imagen}[2]{
	\begin{figure}[!h]
		\centering
		\includegraphics[width=0.9\textwidth]{#1}
		\caption{#2}\label{fig:#1}
	\end{figure}
	\FloatBarrier
}

\newcommand{\imagenflotante}[2]{
	\begin{figure}%[!h]
		\centering
		\includegraphics[width=0.9\textwidth]{#1}
		\caption{#2}\label{fig:#1}
	\end{figure}
}



% El comando \figura nos permite insertar figuras comodamente, y utilizando
% siempre el mismo formato. Los parametros son:
% 1 -> Porcentaje del ancho de página que ocupará la figura (de 0 a 1)
% 2 --> Fichero de la imagen
% 3 --> Texto a pie de imagen
% 4 --> Etiqueta (label) para referencias
% 5 --> Opciones que queramos pasarle al \includegraphics
% 6 --> Opciones de posicionamiento a pasarle a \begin{figure}
\newcommand{\figuraConPosicion}[6]{%
  \setlength{\anchoFloat}{#1\textwidth}%
  \addtolength{\anchoFloat}{-4\fboxsep}%
  \setlength{\anchoFigura}{\anchoFloat}%
  \begin{figure}[#6]
    \begin{center}%
      \Ovalbox{%
        \begin{minipage}{\anchoFloat}%
          \begin{center}%
            \includegraphics[width=\anchoFigura,#5]{#2}%
            \caption{#3}%
            \label{#4}%
          \end{center}%
        \end{minipage}
      }%
    \end{center}%
  \end{figure}%
}

%
% Comando para incluir imágenes en formato apaisado (sin marco).
\newcommand{\figuraApaisadaSinMarco}[5]{%
  \begin{figure}%
    \begin{center}%
    \includegraphics[angle=90,height=#1\textheight,#5]{#2}%
    \caption{#3}%
    \label{#4}%
    \end{center}%
  \end{figure}%
}
% Para las tablas
\newcommand{\otoprule}{\midrule [\heavyrulewidth]}
%
% Nuevo comando para tablas pequeñas (menos de una página).
\newcommand{\tablaSmall}[5]{%
 \begin{table}
  \begin{center}
   \rowcolors {2}{gray!35}{}
   \begin{tabular}{#2}
    \toprule
    #4
    \otoprule
    #5
    \bottomrule
   \end{tabular}
   \caption{#1}
   \label{tabla:#3}
  \end{center}
 \end{table}
}

%
% Nuevo comando para tablas pequeñas (menos de una página).
\newcommand{\tablaSmallSinColores}[5]{%
 \begin{table}[H]
  \begin{center}
   \begin{tabular}{#2}
    \toprule
    #4
    \otoprule
    #5
    \bottomrule
   \end{tabular}
   \caption{#1}
   \label{tabla:#3}
  \end{center}
 \end{table}
}

\newcommand{\tablaApaisadaSmall}[5]{%
\begin{landscape}
  \begin{table}
   \begin{center}
    \rowcolors {2}{gray!35}{}
    \begin{tabular}{#2}
     \toprule
     #4
     \otoprule
     #5
     \bottomrule
    \end{tabular}
    \caption{#1}
    \label{tabla:#3}
   \end{center}
  \end{table}
\end{landscape}
}

%
% Nuevo comando para tablas grandes con cabecera y filas alternas coloreadas en gris.
\newcommand{\tabla}[6]{%
  \begin{center}
    \tablefirsthead{
      \toprule
      #5
      \otoprule
    }
    \tablehead{
      \multicolumn{#3}{l}{\small\sl continúa desde la página anterior}\\
      \toprule
      #5
      \otoprule
    }
    \tabletail{
      \hline
      \multicolumn{#3}{r}{\small\sl continúa en la página siguiente}\\
    }
    \tablelasttail{
      \hline
    }
    \bottomcaption{#1}
    \rowcolors {2}{gray!35}{}
    \begin{xtabular}{#2}
      #6
      \bottomrule
    \end{xtabular}
    \label{tabla:#4}
  \end{center}
}

%
% Nuevo comando para tablas grandes con cabecera.
\newcommand{\tablaSinColores}[6]{%
  \begin{center}
    \tablefirsthead{
      \toprule
      #5
      \otoprule
    }
    \tablehead{
      \multicolumn{#3}{l}{\small continúa desde la página anterior}\\
      \toprule
      #5
      \otoprule
    }
    \tabletail{
      \hline
      \multicolumn{#3}{r}{\small continúa en la página siguiente}\\
    }
    \tablelasttail{
      \hline
    }
    \bottomcaption{#1}
    \begin{xtabular}{#2}
      #6
      \bottomrule
    \end{xtabular}
    \label{tabla:#4}
  \end{center}
}

%
% Nuevo comando para tablas grandes sin cabecera.
\newcommand{\tablaSinCabecera}[5]{%
  \begin{center}
    \tablefirsthead{
      \toprule
    }
    \tablehead{
      \multicolumn{#3}{l}{\small continúa desde la página anterior}\\
      \hline
    }
    \tabletail{
      \hline
      \multicolumn{#3}{r}{\small continúa en la página siguiente}\\
    }
    \tablelasttail{
      \hline
    }
    \bottomcaption{#1}
  \begin{xtabular}{#2}
    #5
   \bottomrule
  \end{xtabular}
  \label{tabla:#4}
  \end{center}
}



\definecolor{cgoLight}{HTML}{EEEEEE}
\definecolor{cgoExtralight}{HTML}{FFFFFF}

%
% Nuevo comando para tablas grandes sin cabecera.
\newcommand{\tablaSinCabeceraConBandas}[5]{%
  \begin{center}
    \tablefirsthead{
      \toprule
    }
    \tablehead{
      \multicolumn{#3}{l}{\small\sl continúa desde la página anterior}\\
      \hline
    }
    \tabletail{
      \hline
      \multicolumn{#3}{r}{\small\sl continúa en la página siguiente}\\
    }
    \tablelasttail{
      \hline
    }
    \bottomcaption{#1}
    \rowcolors[]{1}{cgoExtralight}{cgoLight}

  \begin{xtabular}{#2}
    #5
   \bottomrule
  \end{xtabular}
  \label{tabla:#4}
  \end{center}
}


















\graphicspath{ {./img/} }

% Capítulos
\chapterstyle{bianchi}
\newcommand{\capitulo}[2]{
	\setcounter{chapter}{#1}
	\setcounter{section}{0}
	\chapter*{#2}
	\addcontentsline{toc}{chapter}{#2}
	\markboth{#2}{#2}
}

% Apéndices
\renewcommand{\appendixname}{Apéndice}
\renewcommand*\cftappendixname{\appendixname}

\newcommand{\apendice}[1]{
	%\renewcommand{\thechapter}{A}
	\chapter{#1}
}

\renewcommand*\cftappendixname{\appendixname\ }

% Formato de portada
\makeatletter
\usepackage{xcolor}
\newcommand{\tutor}[1]{\def\@tutor{#1}}
\newcommand{\course}[1]{\def\@course{#1}}
\definecolor{cpardoBox}{HTML}{E6E6FF}
\def\maketitle{
  \null
  \thispagestyle{empty}
  % Cabecera ----------------
\noindent\includegraphics[width=\textwidth]{cabecera}\vspace{1cm}%
  \vfill
  % Título proyecto y escudo informática ----------------
  \colorbox{cpardoBox}{%
    \begin{minipage}{.8\textwidth}
      \vspace{.5cm}\Large
      \begin{center}
      \textbf{TFG del Grado en Ingeniería Informática}\vspace{.6cm}\\
      \textbf{\LARGE\@title{}}
      \end{center}
      \vspace{.2cm}
    \end{minipage}

  }%
  \hfill\begin{minipage}{.20\textwidth}
    \includegraphics[width=\textwidth]{escudoInfor}
  \end{minipage}
  \vfill
  % Datos de alumno, curso y tutores ------------------
  \begin{center}%
  {%
    \noindent\LARGE
    Presentado por \@author{}\\ 
    en Universidad de Burgos --- \@date{}\\
    Tutor: \@tutor{}\\
  }%
  \end{center}%
  \null
  \cleardoublepage
  }
\makeatother

\newcommand{\nombre}{Daniel Santidrián Alonso} %%% cambio de comando

% Datos de portada
\title{UBUassistant}
\author{Daniel Santidrián Alonso}
\tutor{Pedro Renedo Fernández}
\date{\today}

\begin{document}

\maketitle


\newpage\null\thispagestyle{empty}\newpage


%%%%%%%%%%%%%%%%%%%%%%%%%%%%%%%%%%%%%%%%%%%%%%%%%%%%%%%%%%%%%%%%%%%%%%%%%%%%%%%%%%%%%%%%
\thispagestyle{empty}


\noindent\includegraphics[width=\textwidth]{cabecera}\vspace{1cm}

\noindent D. Pedro Renedo Fernández, profesor del departamento de Ingeniería Civil, área de lenguajes y sistemas informáticos.

\noindent Expone:

\noindent Que el alumno D. \nombre, con DNI 71310411S, ha realizado el Trabajo final de Grado en Ingeniería Informática titulado UBUassistant. 

\noindent Y que dicho trabajo ha sido realizado por el alumno bajo la dirección del que suscribe, en virtud de lo cual se autoriza su presentación y defensa.

\begin{center} %\large
En Burgos, {\large \today}
\end{center}

\vfill\vfill\vfill

% Author and supervisor
\begin{minipage}{0.45\textwidth}
\begin{flushleft} %\large
Vº. Bº. del Tutor:\\[2cm]
D. Pedro Renedo Fernández
\end{flushleft}
\end{minipage}
\hfill
%\begin{minipage}{0.45\textwidth}
%\begin{flushleft} %\large
%Vº. Bº. del co-tutor:\\[2cm]
%D. nombre co-tutor
%\end{flushleft}
%\end{minipage}
%\hfill

\vfill

% para casos con solo un tutor comentar lo anterior
% y descomentar lo siguiente
%Vº. Bº. del Tutor:\\[2cm]
%D. nombre tutor


\newpage\null\thispagestyle{empty}\newpage




\frontmatter

% Abstract en castellano
\renewcommand*\abstractname{Resumen}
\begin{abstract}
La búsqueda de información específica en una página web puede ser un proceso largo y tedioso, y más aún cuando accedes a una página web por primera vez. Ciertos portales web, normalmente de instituciones públicas, disponen de abundante información que debería estar visible en un primer plano pero que, sin embargo, está encubierto por diversos menús o submenús.

Con estas premisas nace UBUassistant, un asistente virtual capaz de reconocer una entrada de texto por parte de un usuario y otorgar una respuesta a dicha entrada. De esta manera un usuario puede interactuar con un asistente para encontrar un apartado del portal web en lugar de navegar entre menús que pueden ser ambiguos.

Concretamente, UBUassistant es capaz de orientar a un usuario de la página web de la Universidad de Burgos sin tener que navegar sobre el portal, si existe una respuesta concreta ésta será mostrada de manera inmediata, de lo contrario se otorgarán sugerencias para localizar la repuesta deseada. Basta con preguntar al asistente sobre algún apartado del sitio web.

\end{abstract}

\renewcommand*\abstractname{Descriptores}
\begin{abstract}
Asistente virtual, orientación en portales web, razonamiento basado en casos, aprendizaje tutelado, aplicación web.
\end{abstract}

\clearpage

% Abstract en inglés
\renewcommand*\abstractname{Abstract}
\begin{abstract}
Searching for specific information in web sites can be a long and tedious process, and even more when you acess for first time a web page. Certain web portals, usually public institutions, provide abundant information that should be visible at first sight but, however, is covered by various menus or sub-menus.

With these premises UBUassistant is born, a virtual assistant that is able to recognize a text entry by a user and granting a response to that entry. This way a user can interact with an assistant to find a section of the web inestead of naviganting among menus that may be ambiguous.

Specifically, UBUassistant is able to guide a user of the website of the University of Burgos without having to navigate through the web, if there is a specific response, it will be shown immediately, otherwise suggestions will be given to locate the desired response. Just ask the assistant about some section of the website.
\end{abstract}

\renewcommand*\abstractname{Keywords}
\begin{abstract}
Virtual assistant, orientation in web portals, case based reasoning, tutored learning, web application.
\end{abstract}

\clearpage

% Indices
\tableofcontents

\clearpage

\listoffigures

\clearpage

\listoftables
\clearpage

\mainmatter
\capitulo{1}{Introducción}

Para ciertas personas, poder encontrar información desde un \textit{smartphone} les supone una tarea bastante tediosa y complicada, lo cuál les lleva muchas veces a la deseperación y finalmente no acaban encontrando aquello que deseaban encontrar.

La página web de la Universidad de Burgos es una de las mejores adaptadas a los diferentes dispositivos con los que se puede acceder a la misma. Pero como la gran mayoría de páginas web de universidades, esta contiene mucha información, como es lógico, haciendo que para ciertas personas pueda ser algo complicado navegar a través de ella. Esto se debe que a pesar de que la página esté estructurada en cinco menús, la cantidad de submenús que se encuentran dentro de ellos pueden llegar a confundir al usuario cuando a accedido a varios de ellos desde estos dispositivos, haciéndole dudar como ha conseguido entrar hasta el mismo para futuras consultas.

El método propuesto para facilitar la búsqueda de estos usuarios es una aplicación para dispositivos Android, ya que en España el 87,1 \% de los dispositivos son Android \cite{android:data}. De esta forma se pretende facilitar el uso de la web de la Universidad de Burgos al mayor porcentaje posible de personas.

El asistente es el encargado de buscar una respuesta al texto introducido por el usuario dentro de la aplicación. El usuario no tiene que introducir palabras clave para que el asistente funcione correctamente, sino que puede emplear lenguaje natural, ya que los algoritmos del asistente es el encargado de analizar el texto introducido y poder encontrar así la respuesta más adecuada para el usuario.

\section{Estructura de la memoria}\label{estructura-de-la-memoria}
La memoria tiene la siguiente división en apartados:

\begin{itemize}
	\tightlist
	\item
	\textbf{Introducción:} en este apartado se realiza una descripción de una manera breve del problema que se intenta resolver y la solución otorgada. Además incluye subapartados con la estructura de la memoria y el listado de materiales adjuntos.
	\item
	\textbf{Objetivos del proyecto:} sección donde se explican los objetivos de desarrollar un proyecto de estas características.
	\item
	\textbf{Conceptos teóricos:} capítulo en el que se abordan los conceptos teóricos necesarios para comprender el resultado final del proyecto.
	\item
	\textbf{Técnicas y herramientas:} en esta sección se describen las herramientas y las técnicas que se han utilizado para el desarrollo y gestión del proceso del proyecto.
	\item
	\textbf{Aspectos relevantes del desarrollo:} apartado donde se tratan aquellos aspectos que se consideran destacados en el desarrollo del proyecto.
	\item
	\textbf{Trabajos relacionados:} capitulo que expone y describe aquellos trabajos que están relacionados con la temática de asistente virtual.
	\item
	\textbf{Conclusiones y líneas de trabajo futuras:} sección que explica las conclusiones obtenidas tras la realización del proyecto y la funcionalidad que es posible añadir en el futuro.
\end{itemize}

Además, se proporcionan los siguientes anexos:

\begin{itemize}
	\tightlist
	\item
	\textbf{Plan del proyecto software:} capítulo donde se expone planificación temporal del proyecto y su viabilidad.
	\item
	\textbf{Especificación de requisitos:} en este apartado se desarrollan los objetivos del software y la especificación de requisitos.
	\item
	\textbf{Especificación de diseño:} sección que describe el diseño de datos, el diseño procedimental y el diseño arquitectónico.
	\item
	\textbf{Documentación técnica de programación:} en este capítulo se explica todo lo relacionado con la programación, la estructura de directorios, el manual del programador y las pruebas realizadas.
	\item
	\textbf{Documentación de usuario:} apartado que realiza un explicación sobre los requisitos de usuarios, la instalación y proporciona un manual de usuario.
\end{itemize}

\section{Materiales adjuntos}\label{materiales-adjuntos}

Los materiales que se adjuntan con la memoria son: 

\begin{itemize}
	\tightlist
	\item
	Aplicación Java servlet UBUassistant.
	\item
	Aplicación Android cliente UBUassistant.
	\item
	Máquina virtual configurada para acceder al servidor de forma local.
	\item
	Captura de la máquina virtual de Microsoft Azure usada como servidor global de la aplicación.
	\item	
	JavaDoc.
\end{itemize}

Además, los siguientes recursos están accesibles a través de internet:

\begin{itemize}
	\tightlist
	\item
	Repositorio del proyecto. \cite{ubuassistant:repo}
\end{itemize}

\capitulo{2}{Objetivos del proyecto}

En este apartado se explica cuáles han sido los objetivos que se han perseguido durante el proyecto.

\section{Objetivos software}\label{objetivos-software}

\begin{itemize}
	\tightlist
	\item
	Desarrollar una aplicación Android que permita a cualquier usuario poder acceder de forma sencilla rápida a la diferente información que contiene la página de la Universidad de Burgos utilizando su teléfono móvil.
	\item
	Mejorar la experiencia de usuario a través de su dispositivo móvil.
	\item
	Realizar un aprendizaje de nuevos casos en los que no existe respuesta de manera supervisada.
\end{itemize}

\section{Objetivos técnicos}\label{objetivos-tecnicos}

\begin{itemize}
	\tightlist
	\item
	Aplicar Scrum, en la medida de lo posible, como metodología de desarrollo ágil.
	\item
	Realizar la aplicación siguiendo el concepto de Modelo Vista Controlador, separando la interfaz de usuario, el motor de la aplicación y los datos.
	\item
	Obtener las respuestas del motor de la aplicación empleando un estándar, haciendo así que la aplicación sea más modular y multiplataforma.
	\item
	Acceder a una base de datos MySQL mediante Hibernate y JDBC.
	\item
	Servirse de GitHub como sistema de control de versiones.
	\item
	Utilizar herramientas de control de calidad del software como SonarQube, RefactorIt o InCode.
	
\end{itemize}

\section{Objetivos personales}\label{objetivos-personales}

\begin{itemize}
	\tightlist
	\item
	Emplear la mayor cantidad posible de conocimientos adquiridos durante la carrera.
	\item
	Ampliar los conocimientos adquiridos durante la carrera utilizando nuevos lenguajes de programación (Android).
	\item 
	Conseguir una aplicación que facilite las tareas de búsqueda en la página de la Universidad de Burgos.
\end{itemize}

\capitulo{3}{Conceptos teóricos}

En aquellos proyectos que necesiten para su comprensión y desarrollo de unos conceptos teóricos de una determinada materia o de un determinado dominio de conocimiento, debe existir un apartado que sintetice dichos conceptos.

Algunos conceptos teóricos de \LaTeX \footnote{Créditos a los proyectos de Álvaro López Cantero: Configurador de Presupuestos y Roberto Izquierdo Amo: PLQuiz}.

\section{Secciones}

Las secciones se incluyen con el comando section.

\subsection{Subsecciones}

Además de secciones tenemos subsecciones.

\subsubsection{Subsubsecciones}

Y subsecciones. 


\section{Referencias}

Las referencias se incluyen en el texto usando cite \cite{wiki:latex}. Para citar webs, artículos o libros \cite{koza92}.


\section{Imágenes}

Se pueden incluir imágenes con los comandos standard de \LaTeX, pero esta plantilla dispone de comandos propios como por ejemplo el siguiente:

\imagen{escudoInfor}{Autómata para una expresión vacía}



\section{Listas de items}

Existen tres posibilidades:

\begin{itemize}
	\item primer item.
	\item segundo item.
\end{itemize}

\begin{enumerate}
	\item primer item.
	\item segundo item.
\end{enumerate}

\begin{description}
	\item[Primer item] más información sobre el primer item.
	\item[Segundo item] más información sobre el segundo item.
\end{description}
	
\begin{itemize}
\item 
\end{itemize}

\section{Tablas}

Igualmente se pueden usar los comandos específicos de \LaTeX o bien usar alguno de los comandos de la plantilla.

\tablaSmall{Herramientas y tecnologías utilizadas en cada parte del proyecto}{l c c c c}{herramientasportipodeuso}
{ \multicolumn{1}{l}{Herramientas} & App AngularJS & API REST & BD & Memoria \\}{ 
HTML5 & X & & &\\
CSS3 & X & & &\\
BOOTSTRAP & X & & &\\
JavaScript & X & & &\\
AngularJS & X & & &\\
Bower & X & & &\\
PHP & & X & &\\
Karma + Jasmine & X & & &\\
Slim framework & & X & &\\
Idiorm & & X & &\\
Composer & & X & &\\
JSON & X & X & &\\
PhpStorm & X & X & &\\
MySQL & & & X &\\
PhpMyAdmin & & & X &\\
Git + BitBucket & X & X & X & X\\
Mik\TeX{} & & & & X\\
\TeX{}Maker & & & & X\\
Astah & & & & X\\
Balsamiq Mockups & X & & &\\
VersionOne & X & X & X & X\\
} 

\capitulo{4}{Técnicas y herramientas}

\section{Patrones de diseño}
\subsection{Modelo Vista Controlador}

Modelo vista controlador (\emph{MVC}) es un patrón arquitectónico que nos ayuda a separar los datos, la lógica de negocio y la interfaz de usuario \cite{mvc:wiki}.

\imagen{MVC}{Diagrama Modelo Vista Controlador}

El significado de cada uno de los tres componentes de este modelo son:

\begin{itemize}
	\tightlist
	\item
	\textbf{Modelo:} representa los datos que utiliza la aplicación.
	\item
	\textbf{Vista:} Muestra al usuario la información.
	\item
	\textbf{Controlador:} Controla las interacciones con el usuario.
\end{itemize}


Estos conceptos pueden ser aplicados a nuestra aplicación Android, obteniendo así un Modelo Vista Controlador muy bien diferenciado y estructurado. Esto lo podemos ver claramente en el siguiente diagrama.
\imagen{MVC-Android}{Modelo Vista Controlador en Android}

\section{Control de versiones}

\begin{itemize}
	\tightlist
	\item
	Herramientas consideradas: \href{https://git-scm.com/}{Git}.
	\item
	Herramienta elegida: \href{https://git-scm.com/}{Git}.
\end{itemize}

Git es un software de control de versiones pensado para proyectos que poseen una gran cantidad de ficheros fuente cuyo propósito es registrar todos los cambios efectuados en dichos ficheros. Además es un software de código libre distribuido bajo licencia \textit{GPL GNU} \cite{git:wiki}.

\section{Hosting del repositorio}

\begin{itemize}
	\tightlist
	\item
	Plataformas consideradas: \href{https://bitbucket.org/}{Bitbucket} y \href{https://github.com/}{Github}.
	\item
	Plataforma elegida: \href{https://github.com/}{Github}.
\end{itemize}

GitHub es una plataforma que es utilizada para alojar proyectos, los cuales emplean Git como sistema de control de versiones.

GitHub ha sido elegido frente a BitBucket debido a que es una plataforma que se ha ido utilizando en diferentes asignaturas como Gestión de Proyectos.

\section{Gestión del repositorio}

\begin{itemize}
	\tightlist
	\item
	Herramientas consideradas: \href{https://www.gitkraken.com/}{GitKraken} y \href{https://desktop.github.com/}{GitDesktop}.
	\item
	Herramienta elegida: \href{https://www.gitkraken.com/}{GitKraken}.
\end{itemize}

GitKraken es una aplicación que sirve para gestionar de una forma más sencilla nuestro repositorio de GitHub. Es una herramienta multiplataforma compatible con Windows, Mac y Linux.

La decisión de usar GitKraken frente a GitDesktop está basada en la experiencia personal, ya que previamente había trabajado con ambas herramientas. A parte de esto, la decisión de usar GitKraken está fundamentada en su compatibilidad con sistemas Linux.

\section{Sistema Operativo}

\begin{itemize}
	\tightlist
	\item
	Sistemas considerados: \href{https://www.microsoft.com/es-es/windows}{Windows 10}, \href{https://www.ubuntu.com/}{Ubuntu} y \href{https://linuxmint.com/}{Linux Mint}.
	\item
	Sistema elegido: \href{https://linuxmint.com/}{Linux Mint}.
\end{itemize}

Linux Mint es un sistema operativo Linux que utiliza un núcleo de sistema basado en Debian y Ubuntu.

En un primer momento se planteó si usar Windows o Linux como sistema operativo. Finalmente se decidió usar Linux ya que se consideró que la gestión de recursos la hace de forma más eficiente y emplea menos recursos. Una vez tomada la decisión de usar Linux, se barajaron dos opciones, si usar Ubuntu o Linux Mint. La decisión se tomo en base a la experiencia personal, ya que se había trabajado con anterioridad con ambos sistemas operativos.

\section{Entorno de desarrollo integrado (IDE) para JAVA}

\begin{itemize}
	\tightlist
	\item
	IDE's considerados: \href{https://www.jetbrains.com/idea/}{IntelliJ}, \href{https://netbeans.org/}{NetBeans} y \href{https://eclipse.org/}{Eclipse}.
	\item
	IDE elegidos: \href{https://eclipse.org/}{Eclipse}.
\end{itemize}

Eclipse es un software que está formado por diferentes herramientas de programación de código abierto multiplataforma y se encuentra bajo la licencia de software libre \textit{Eclipse Public License}. Para el desarrollo de este proceso, se ha empleado la versión EE, la cual permite realizar páginas web basadas en Java en formato JSP \cite{eclipse:info}.

De todos los softwares valorados para emplear como IDE se decidió utilizar Eclipse ya que es una herramienta que resulta familiar debido a que ha sido utilizada en todas aquellas asignaturas del Grado en el que se ha utilizado el lenguaje de programación JAVA.


\section{Entorno de desarrollo integrado (IDE) para Android}

\begin{itemize}
	\tightlist
	\item
	IDE's considerados: \href{https://developer.android.com/studio/}{Android Studio}, \href{http://www.android-ide.com/}{AIDE} y \href{https://eclipse.org/}{Eclipse}.
	\item
	IDE elegidos: \href{https://developer.android.com/studio/}{Android Studio}.
\end{itemize}

Android Studio es el IDE oficial para el desarrollo de aplicaciones para Android. Este IDE es el sucesor de Eclipse como IDE oficial de desarrollo, funciona gratuitamente bajo \textit{Licencia Apache 2.0} y está disponible para Windows, macOS y Linux \cite{androidstudio:wiki}.

Esta herramienta ha sido elegida principalmente por que es la herramienta oficial que nos ofrece Google y trae consigo todo lo necesario para desarrollar la aplicación, sin tener que descargar complementos para el software, como pasa en el caso de Eclipse.

\section{Documentación}

\begin{itemize}
	\tightlist
	\item
	Herramientas consideradas: \href{http://www.xm1math.net/texmaker/}{TexMaker}, \href{https://www.texstudio.org/}{TexStudio}, \href{https://es.sharelatex.com/}{ShareLaTex} y \href{https://www.openoffice.org/es/}{OpenOffice}.
	\item
	Herramienta elegida: \href{http://www.xm1math.net/texmaker/}{TexMaker}.
\end{itemize}

Entre realizar la documentación con LaTex u OpenOffice se decidió emplear LaTex, ya que los resultados que se obtienen son una documentación obtenida con una gran calidad tipográfica.

LaTex es un sistema empleado para la composición de textos escritos con una alta calidad tipográfica tales como artículos académicos, tesis y libros técnicos. Es un software libre bajo la licencia \textit{LPPL} \cite{latex:wiki}.

TexStudio es un editor para LaTex de código abierto y multiplataforma que es una evolución del IDE TexMaker. Este IDE esta bajo la licencia \textit{GNU GPL} \cite{texstudio:wiki}.

La elección de este editor es que es offline, por lo que no dependemos de una conexión a Internet, y que es una evolución de TexMaker, por lo que funcionalidades añadidas respecto a este IDE.

\capitulo{5}{Aspectos relevantes del desarrollo del proyecto}

Este apartado pretende recoger los aspectos más interesantes del desarrollo del proyecto, comentados por los autores del mismo. Aquí se incluyen las diferentes explicaciones tanto del diseño como de la implementación.

\section{Inicio y fase de análisis}
Este proyecto es la consecución del TFG UBUassistant de Daniel Santidrián Alonso en el cual durante la defensa del mismo, surgieron diferentes mejoras para el mismo \cite{tfg:art}.

Tras valorar las diferentes mejoras con el tutor del proyecto, se decidió continuar con él, debido a que me pareció un proyecto muy interesante tanto a nivel de usuario como a nivel de programación.

\section{Metodologías aplicadas}
En el desarrollo de este proyecto se ha intentado emplear en la medida de los posible la metodología ágil de Scrum.

Debido a que el tamaño del proyecto es pequeño, no se ha podido seguir de forma estricta todas las pautas de la metodología, como las reuniones diarias con todo el equipo. Las pautas que se han seguido durante el proyecto han sido:
\begin{itemize}
	\tightlist
	\item
	Desarrollo incremental del proyecto mediante sprints.
	\item
	Spints de duración semanal en vez de diaria por el tamaño del proyecto.
	\item
	Al finalizar el sprint, reuniones para evaluar el proyecto y plantear los pasos a seguir en el siguiente sprint.
\end{itemize}

Para el desarrollo del servidor del asistente, se realizaron diferentes pruebas de ensayo error hasta que se obtuvieron los datos a través del estandar JSON de la forma deseada.

Al final de cada sesión de trabajo se ha intentado realizar un \textit{commit} en el repositorio del proyecto para mantener así un control sobre los diferentes pasos que se han ido dando.

\section{Formación}

Durante las primeras fases del proyecto, fue necesario aprender el funcionamiento del TFG de Daniel Santidrián Arce. Para esto, se consultó la documentación de dicho proyecto, así como el código fuente del mismo \cite{tfg:art}.

A parte de esto, se realizó una investigación sobre el funcionamiento genérico de un Sistema de Razonamiento Basado en Casos (\textit{CBR}) como de la utilidad empleada en la versión anterior del TFG, \textit{jCOLIBRI}. Para ello se consultaron los mismos artículos empleados para este propósito por el autor del anterior TFG.

\begin{itemize}
	\tightlist
	\item
	Razonamiento Basado en Casos: Una visión general (Laura Lozano y Javier Fernández) \cite{cbr:art}.
	\item 
	Tutorial jCOLIBRI (J. Recio García, B. Díaz Aguado y P. González Calero) \cite{colibri:tut}.
\end{itemize}

\capitulo{6}{Trabajos relacionados}

Existen diferentes métodos de implementar un asistente virtual que permita el reconocimiento de texto. La amplia mayoría utilizan algoritmos de inteligencia artificial, siendo algunos de ellos los que se decantan por el razonamiento basado en casos.

\section{Artículos de investigación}

\subsection{Comportamiento Adaptable de Chatbots Dependiente del Contexto}

Artículo desarrollado por el grupo de investigación en sistemas de
Información (GISI) de la Universidad Nacional de Lanús en Argentina, en el que se estudian los diferentes algoritmos utilizados en asistentes virtuales o \emph{chatbots} centrándose en las mejoras que pueden introducirse en el reconocimiento del contexto. En el artículo se plantea el razonamiento basado en casos como una de las soluciones \cite{rev:invest}.

\subsection{Un Sofbot de Charla Desarrollado con Técnicas de Razonamiento basado en Casos}

En este artículo de la Revista de Investigación de Sistemas e Informática de la Universidad Nacional Mayor de San Marcos se exponen las técnicas utilizadas para la construcción de un software de charla. El software construido utiliza la técnica de razonamiento basado en casos y se explican los distintos pasos realizados para su creación \cite{SAMI:project}.


\section{Proyectos}

\subsection{ChatBot}

Es un proyecto \emph{Open Source} que desarrolla una aplicación web donde se hospeda un chat construido mediante razonamiento basado en casos. Utiliza el framework de jCOLIBRI para llevarlo a cabo. Es un pequeño proyecto desarrollado por un estudiante de la Universidad de Auckland \cite{ChatBot:project}.

\begin{itemize}
\tightlist
\item
Web del proyecto: https://github.com/agkphysics/ChatBot
\end{itemize}

\subsection{SAMI}

En el anterior artículo de investigación \cite{SAMI:project}, se explican las técnicas seguidas para la realización de un software de charla basado en razonamiento basado en casos. El finalidad de este software consiste en proporcionar a los alumnos del aula virtual de la universidad información acerca de la misma. El software reconoce el diálogo escrito.

Este proyecto no dispone de página web y no ha sido encontrada más información a parte de la documentación.

\section{Fortalezas y debilidades del proyecto}

Las fortalezas de este proyecto radican en:

\begin{itemize}
\tightlist
\item
La interfaz de usuario es muy sencilla.
\item
La interfaz es completa y se permite interactuar al usuario con ella.
\item
La aplicación es adaptable a cualquier ámbito solo con cambiar los casos de la base de datos.
\item
En cualquier momento es posible cambiar o retocar los casos, tanto descriptivos como respuesta.
\item
La aplicación dispone de una página de administración para consultar el log de uso y sugerencias de aprendizaje así como la inclusión, la modificación y el borrado de casos.
\item
El asistente virtual puede incluirse en cualquier página web.
\item
La interfaz de la aplicación es adaptable a distintos tamaños de pantalla.
\end{itemize}

Por contra, existen ciertas debilidades como:

\begin{itemize}
\tightlist
\item
La aplicación no está hospedada en ningún sitio web, sólo es accesible mediante una conexión local quedando como una propuesta de futuro poder integrarse en algún servicio de la Universidad de Burgos.
\item
Actualmente solo está preparada para trabajar con casos que tratan sobre la Universidad de Burgos.
\item
No es todo lo precisa que un usuario podría desear, cubre los apartados más amplios de la Universidad.
\end{itemize}


\capitulo{7}{Conclusiones y Líneas de trabajo futuras}

\section{Conclusiones}

Tras el desarrollo del proyecto, las conclusiones a las que se ha llegado son las siguientes.

\begin{itemize}
	\tightlist
	\item
	Finalmente se ha conseguido realizar el proyecto de la forma en la que se había planteado desde un primer momento, cumpliendo todos los objetivos que se habían marcado al comienzo del proyecto.
	\item 
	
	\item 
	Se han utilizado la gran mayoría de los conocimientos adquiridos a lo largo de la carrera.
	\item 
	Durante el desarrollo del proyecto se han ampliado conocimientos, tales como desarrollo de aplicaciones Android, conexiones cliente-servidor, inteligencia artificial, etc.
	\item 
	El proyecto ha sido posible completarlo dentro del periodo estipulado para ello.
\end{itemize}

\section{Líneas de trabajo futuras}

Este proyecto tiene todavía mucho potencial para crecer aún más, por lo que se puede seguir trabajando en las siguientes líneas:

\begin{itemize}
	\tightlist
	\item
	Crear una aplicación de administración para dispositivos móviles.
	\item 
	Permitir que el algoritmo recoja ámbitos más concretos de la página de la Universidad de Burgos.
	\item
	Realizar otros clientes para otro tipos de dispositivos.
	\item 
	Alojar el servidor de la aplicación en un servidor propiedad de la Universidad de Burgos.
\end{itemize}


\bibliographystyle{plainnat}
\bibliography{bibliografia}

\end{document}
